% GENERATED FILE. Edit canonical lessons, then run npm run generate:books.
% canonical-source-hash: fnv1a64:f9d24b793ae66382
% canonical-lessons: UR-C185-goli, UR-C185-nuskha, UR-C185-injekshan, UR-C185-zakhm, UR-C185-chot

\chapter{Pill, Prescription, Injection, Wound, Injury}
\label{ch:ur-pill-prescription-injection-wound-injury}
\hlchaptermodality{185}

\begin{chapteropening}
\textbf{By the end of this chapter:} \emph{I can say a pill, a prescription, an injection, a wound and an injury.}

Complete the last lesson of chapter 185: I can say a pill, a prescription, an injection, a wound and an injury.
\end{chapteropening}

\section[\texorpdfstring{\ur{گولی} (golī)}{گولی (golī)}]{\ur{گولی} (golī) — a pill; a bullet}

\label{lesson:UR-C185-goli}



\begin{quote}
Before the new one: say the Urdu for to get wet, then the Urdu for to travel.
\end{quote}

\subsection*{You'll want to know: \ur{گولی}}
\textbf{\ur{گولی}} — \emph{golī} — \textquotedblleft{}a pill; a bullet\textquotedblright{}.

\begin{cousinweb}
\textbf{\ur{گولی}} (\emph{golī}) is a tadbhava, a word inherited from Sanskrit and reshaped by sound change. It comes from \emph{golikā}, \textquotedblleft{}little ball\textquotedblright{} (CDIAL 4321). It is built on \emph{gola}, \textquotedblleft{}ball\textquotedblright{}.
\end{cousinweb}

\begin{grammarlens}[title={What you've built}]
One more word for this chapter.
\end{grammarlens}

\subsection*{Your turn}
\begin{itemize}
\raggedright
  \item \emph{Say it:} \emph{golī}
  \item \emph{Say it:} \emph{golī}, once more
  \item \emph{Say it:} say \emph{golī}
  \item \emph{Recall:} say the Urdu for to get wet, then the Urdu for to travel, then say \emph{golī} again
\end{itemize}

\begin{tcolorbox}[breakable,colback=teal!4,colframe=teal!35!black,title={Before you move on}]
What does \ur{گولی} mean? (\textquotedblleft{}A pill; a bullet\textquotedblright{}.) Say it once more.
\end{tcolorbox}
\section[\texorpdfstring{\ur{نسخہ} (nuskhā)}{نسخہ (nuskhā)}]{\ur{نسخہ} (nuskhā) — a prescription}

\label{lesson:UR-C185-nuskha}



\begin{quote}
Before the new one: say the Urdu for to travel, then the Urdu for a pill; a bullet.
\end{quote}

\subsection*{You'll want to know: \ur{نسخہ}}
\textbf{\ur{نسخہ}} — \emph{nuskhā} — \textquotedblleft{}a prescription\textquotedblright{}.

\begin{cousinweb}
\textbf{\ur{نسخہ}} (\emph{nuskhā}) is a loan from Arabic \emph{nuskha}, \textquotedblleft{}a copy, a written text\textquotedblright{}, which reached Urdu through Persian. Its Arabic root is \emph{n-s-kh} (Lane). Urdu uses it for a doctor's written prescription.
\end{cousinweb}

\begin{grammarlens}[title={What you've built}]
One more word for this chapter.
\end{grammarlens}

\subsection*{Your turn}
\begin{itemize}
\raggedright
  \item \emph{Say it:} \emph{nuskhā}
  \item \emph{Say it:} \emph{nuskhā}, once more
  \item \emph{Say it:} say \emph{nuskhā}
  \item \emph{Recall:} say the Urdu for to travel, then the Urdu for a pill; a bullet, then say \emph{nuskhā} again
\end{itemize}

\begin{tcolorbox}[breakable,colback=teal!4,colframe=teal!35!black,title={Before you move on}]
What does \ur{نسخہ} mean? (\textquotedblleft{}A prescription\textquotedblright{}.) Say it once more.
\end{tcolorbox}
\section[\texorpdfstring{\ur{انجکشن} (injekshan)}{انجکشن (injekshan)}]{\ur{انجکشن} (injekshan) — an injection}

\label{lesson:UR-C185-injekshan}



\begin{quote}
Before the new one: say the Urdu for a pill; a bullet, then the Urdu for a prescription.
\end{quote}

\subsection*{You'll want to know: \ur{انجکشن}}
\textbf{\ur{انجکشن}} — \emph{injekshan} — \textquotedblleft{}an injection\textquotedblright{}.

\begin{cousinweb}
\textbf{\ur{انجکشن}} (\emph{injekshan}) is a loan from English \emph{injection}.
\end{cousinweb}

\begin{grammarlens}[title={What you've built}]
One more word for this chapter.
\end{grammarlens}

\subsection*{Your turn}
\begin{itemize}
\raggedright
  \item \emph{Say it:} \emph{injekshan}
  \item \emph{Say it:} \emph{injekshan}, once more
  \item \emph{Say it:} say \emph{injekshan}
  \item \emph{Recall:} say the Urdu for a pill; a bullet, then the Urdu for a prescription, then say \emph{injekshan} again
\end{itemize}

\begin{tcolorbox}[breakable,colback=teal!4,colframe=teal!35!black,title={Before you move on}]
What does \ur{انجکشن} mean? (\textquotedblleft{}An injection\textquotedblright{}.) Say it once more.
\end{tcolorbox}
\section[\texorpdfstring{\ur{زخم} (zakhm)}{زخم (zakhm)}]{\ur{زخم} (zakhm) — a wound}

\label{lesson:UR-C185-zakhm}



\begin{quote}
Before the new one: say the Urdu for a prescription, then the Urdu for an injection.
\end{quote}

\subsection*{You'll want to know: \ur{زخم}}
\textbf{\ur{زخم}} — \emph{zakhm} — \textquotedblleft{}a wound\textquotedblright{}.

\begin{cousinweb}
\textbf{\ur{زخم}} (\emph{zakhm}) is a loan, usually traced to Persian \emph{zakhm}, \textquotedblleft{}a wound\textquotedblright{}.
\end{cousinweb}

\begin{grammarlens}[title={What you've built}]
One more word for this chapter.
\end{grammarlens}

\subsection*{Your turn}
\begin{itemize}
\raggedright
  \item \emph{Say it:} \emph{zakhm}
  \item \emph{Say it:} \emph{zakhm}, once more
  \item \emph{Say it:} say \emph{zakhm}
  \item \emph{Recall:} say the Urdu for a prescription, then the Urdu for an injection, then say \emph{zakhm} again
\end{itemize}

\begin{tcolorbox}[breakable,colback=teal!4,colframe=teal!35!black,title={Before you move on}]
What does \ur{زخم} mean? (\textquotedblleft{}A wound\textquotedblright{}.) Say it once more.
\end{tcolorbox}
\section[\texorpdfstring{\ur{چوٹ} (choṭ)}{چوٹ (choṭ)}]{\ur{چوٹ} (choṭ) — an injury, a knock}

\label{lesson:UR-C185-chot}



\begin{quote}
Before the new one: say the Urdu for an injection, then the Urdu for a wound.
\end{quote}

\subsection*{You'll want to know: \ur{چوٹ}}
\textbf{\ur{چوٹ}} — \emph{choṭ} — \textquotedblleft{}an injury, a knock\textquotedblright{}.

\begin{cousinweb}
\textbf{\ur{چوٹ}} (\emph{choṭ}) is a tadbhava, a word inherited and reshaped by sound change. CDIAL lists Hindi \emph{coṭ}, \textquotedblleft{}wound\textquotedblright{}, under Sanskrit \emph{cuṭati}, \textquotedblleft{}cuts off, strikes\textquotedblright{}, in its form \emph{*coṭṭ-}, and calls the whole group non-Aryan in origin (CDIAL 4857).
\end{cousinweb}

\begin{grammarlens}[title={What you've built}]
One more word for this chapter.
\end{grammarlens}

\subsection*{Your turn}
\begin{itemize}
\raggedright
  \item \emph{Say it:} \emph{choṭ}
  \item \emph{Say it:} \emph{choṭ}, once more
  \item \emph{Say it:} say \emph{choṭ}
  \item \emph{Recall:} say the Urdu for an injection, then the Urdu for a wound, then say \emph{choṭ} again
\end{itemize}

\begin{tcolorbox}[breakable,colback=teal!4,colframe=teal!35!black,title={Before you move on}]
What does \ur{چوٹ} mean? (\textquotedblleft{}An injury, a knock\textquotedblright{}.) Say it once more.
\end{tcolorbox}
